%-*- ES -*-
%----------------------------------------------------------------------
% capitulo 1: Introduccion
%----------------------------------------------------------------------
% Estructura del capitulo:
%
%----------------------------------------------------------------------


Lorem ipsum dolor sit amet, consectetur \cite{Garcia:1997} adipiscing elit, sed do eiusmod tempor incididunt ut labore et dolore magna aliqua. Ut enim ad minim veniam, quis nostrud exercitation ullamco laboris nisi ut aliquip ex ea commodo consequat. Duis aute irure dolor in reprehenderit in voluptate velit esse cillum dolore eu fugiat nulla pariatur. Excepteur sint occaecat cupidatat non proident, sunt in culpa qui officia deserunt mollit anim id est laborum

% se pone una introducci�n al tema en general

\section{Objetivos}

\begin{itemize}
 \item item 1
 \item item 2
\end{itemize}


\begin{verbatim}
 Este es un bien ambiente para 
poner codigo
\end{verbatim}



En \cite{Davis:1989} se \footnote{Esta es una nota al pie} ve...
\\
\textit{casa} \\
\emph{casa} \\
\texttt{casa} \\
\textsf{casa} \\
\textbf{casa} \\

\subsection{objetivos principales}

\subsubsection{particulares}

\begin{center}
\begin{figure}

\includegraphics[width=3cm,height=3cm]{UniLogo}

\caption{Esta es la figura del escudo de la uns}
\end{figure}
\end{center}

\section{Plan de tesis y principales contribuciones}\label{secContribuciones}

\section{Trabajos previos relacionados}
% rese�ar los art�culos hechos en el contexto del trabajo de tesis si los hay

